% !TEX root = main.tex
\subsubsection*{전달 방식의 기능과 비용}
직접 전달은 Receiver Note와 Change Note를 한 transaction에서 만든다.
Voucher 경로는 Transfer 이후 Receiver의 Proceed 또는 Sender의 Recall로 해결된다.
정상 수령과 회수는 최종 소유자가 다르므로 경로를 구분한다.
\begin{center}\small\setlength{\tabcolsep}{6pt}
    \textbf{표 7. 전달·회수 경로의 Gas}\par\smallskip
    \begin{tabular}{@{}lcrl@{}}\toprule
        \textbf{경로}            & \textbf{Audit} & \textbf{Gas} & \textbf{통계·값의 성격} \\ \midrule
        직접 전달 기준선              & 없음             & 2,366,228    & 5회 median         \\
        SharedSecret 직접 전달     & 있음             & 2,724,242    & 공식 1회             \\
        Transfer $\to$ Proceed & 있음             & 5,361,496    & Event별 median 합   \\
        Transfer $\to$ Recall  & 있음             & 5,362,498    & Event별 median 합   \\
        \bottomrule\end{tabular}\end{center}

Proceed 경로는 3,577,414+1,784,082이며, Recall 경로는 3,577,414+1,785,084이다.
두 값은 각각의 대표 Event median을 합한 값이다. 동일 경로를 한 번에 실행해 측정한 Gas median이나 E2E가 아니다.
직접 전달과의 차이에는 Audit, 추가 transaction, 출력·Tree 작업의 차이가 함께 들어가므로 Voucher만의 순수 비용으로 해석하지 않는다.

SharedSecret 대안은 공유 opening에서 파생한 값으로 Receiver Note의 nf를 계산한다.
Sender도 그 값을 알 수 있어 다음 소비를 연결할 수 있지만, 소비에는 별도로 Receiver secret이 필요하다.
새 측정은 실제 감사 자료 저장·복호화와 Receiver Exit의 동일 nf 소비를 확인했다.

Voucher에서는 Sender가 rvnf와 Proceed의 출력을 연결할 수 있지만 Receiver Note의 nf는 Receiver secret에 의존한다.
감사자는 각 기록의 암호문을 통해 소비 식별자를 복원하여 다음 거래를 찾는다.
이 구분은 공개 시각·메타데이터 등 다른 단서에 의한 추론까지 모두 배제한다는 뜻이 아니다.
    {\small 근거: [D1,D5]. SharedSecret nf는 기존 owner-secret nf와 다른 실험용 관계이며, 별도 Ledger를 사용했다.}\par

